% !PRINTED 170
\weberpage{170}

\begin{center}
§. 52.

\gesperrt{The Tschirnhausen Transformation.}
\end{center}

We now make one more important application of the theory of symmetric functions and of the formation of resultants, to the transformation of algebraic equations that goes back to \gesperrt{Tschirnhausen}\textsuperscript{1)}.

Let
\weq{1}{f(x) = x^n + a_1 x^{n-1} + a_2 x^{n-2} + \cdots + a_{n-1} x + a_n}
be any function of degree $n$, and let
\weq{2}{y = \frac{\chi(x)}{\psi(x)}}
be any (integral or fractional) rational function of $x$, subject to the single assumption that $f(x)$ and $\psi(x)$ \gesperrt{are to have no common divisor}. By the concluding theorem of \S.\,6 we can then determine the functions $F(x)$ and $\varphi(x)$, the latter of degree at most $n-1$, so that
\weq{3}{\chi(x) = F(x) f(x) + \varphi(x) \psi(x)}
holds; consequently, as soon as a root of the equation
\weq{4}{f(x) = 0}
is put for $x$, we have
\[
\frac{\chi(x)}{\psi(x)} = \varphi(x). % [sic] period as printed before the following continuation
\]
Since we shall consider the values of $y$ only for those values of $x$ that are roots of (4), we lose nothing in generality if from the outset we set
\[
y = \varphi(x),
\]
that is, equal to an integral function of degree $(n-1)$. Let us therefore put
\weq{5}{y = \varphi(x) = \alpha_0 + \alpha_1 x + \alpha_2 x^2 + \cdots + \alpha_{n-1} x^{n-1},}
in which the coefficients $\alpha$ remain for the present entirely undetermined.

If we denote by
\weq{6}{x_1, x_2 \ldots x_n}
%CONT>>

\par\smallskip
{\small 1) Acta eruditorum, Leipzig 1683.}\par
